% document formatting
\documentclass[10pt]{article}
\usepackage[utf8]{inputenc}
\usepackage[left=1in,right=1in,top=1in,bottom=1in]{geometry}
\usepackage[T1]{fontenc}
\usepackage{xcolor}

% math symbols, etc.
\usepackage{amsmath, amsfonts, amssymb, amsthm}

% lists
\usepackage{enumerate}

% images
\usepackage{graphicx} % for images
\usepackage{tikz}

% code blocks
% \usepackage{minted, listings} 

% verbatim greek
\usepackage{alphabeta}

\graphicspath{{./assets/images/Week 2}}

\title{02-614 Week 2 \\ \large{String Algorithms}}
 
\author{Aidan Jan}

\date{\today}

\begin{document}
\maketitle

\section*{Seminumerical String Algorithms}
\subsection*{Rabin-Karp Fingerprinting}
This is an exact matching algorithm exploiting the fact that characters are stored in a computer as binary.  This makes it a \textit{seminumerical} string algorithm.\\\\
Recall the exact matching problem: text $T$, pattern $P$.  Find all occurrences of $P$ in $T$.\\\\
For Rabin-Karp Fingerprinting, let
\begin{itemize}
	\item $\Sigma = \{\text{``0'', ``1'', \dots, ``9''}\}$.  (Note these are strings, not integers.)
	\item $T = \text{``4278931''}$  (Note this is a string, not a integer.)
	\item $P = \text{``78''}$
\end{itemize}
Now, since $P$ has length 2, consider every substring of $T$ of length 2.  $t_i=$ (``42'', ``27'', \dots, ``31''.)  If we convert all substring in $T$ to integers (so we have an integer array), and we convert $P$ to an integer, $p$, then doing the matching is trivial.  This problem essentially becomes finding when an integer appears in a string.  The issue is that generating that integer array is not trivial.

\subsubsection*{Converting $P$ to $p$}
Let $m = |P|$.  Then, (assume 1-indexed),
\[p = P[m] + 10(P[m - 1] + (10P[m - 2] + \dots + P[1]))\]
This formula is known as \textit{Horner's rule}.
\begin{itemize}
	\item Note that none of the numbers in this formula is larger than the size of our alphabet.
	\item Each character in the string using this formula is $O(1)$ work, since it is a single multiplication in the recursive formula.
	\item Therefore, $P \rightarrow p$ runs in $O(|P|)$ time.
\end{itemize}

\subsubsection*{Converting $T$ to $t$}
Let
\[t_i = 10 \times (t_{i - 1} - 10^{m - 1} T[i - 1]) \times T[s + m - 1]\]
This formula is similar to Horner's rule; it essentially does:
\begin{itemize}
	\item Calculate the first integer substring, $t_1$
	\item Subtract the leading digit of the first integer, multiply by 10 (or size of alphabet), and add the next integer to generate $t_2$
	\item Repeat until the end of the string to calculate all the $t_i$ values.
\end{itemize}
This makes the array of integers in $O(|T|)$ time, since each $t_i$ value occurs in constant time.  Once again, it is a single multiplication per character.
\begin{itemize}
	\item Note that this would be even faster in binary, where the alphabet size is $2$.  In this case, multiplications become bit shifts.
\end{itemize}

\subsection*{Issues With This Method}
The primary issue is that we did not state when computers are good at multiplication.  If $T$ is a billion base-pair genome and $P$ is a thousand base-pair read, then the computer will struggle to multiply anything.\\\\
\textbf{Problem:} $t_i$ and $p$ could be huge numbers.  Instead, we can do the following:
\begin{itemize}
	\item Define $t_i' \equiv t_i \mod 8$
	\item Define $p' \equiv p \mod 8$
	\item $P$ occurs at point $i$ in $T$ implies $t_i' = p'$.  Unfortunately, this implication only goes one direction, and we cannot assume the opposite.
	\item $t_i' = p'$ does \textbf{NOT} imply that $P$ occurs at $i$.  We will now have false positives.
\end{itemize}

\subsubsection*{Dealing with False Positives}
Let $q$ be a prime number less than or equal to $M$.  Let $M$ be a number that we pick such that any $q$ less than $M$ is small enough to fit into our computer's registers.  Now, choose a $q$ randomly.
\begin{itemize}
	\item We deal with the false positive issue by adding some randomness.  So essentially, false positives would only occur if we get really unlucky.
\end{itemize}
\textbf{Theorem:} if $M$ is big enough, then
\[\text{Pr}_q[t_i' = p'] \leq 0.2 \text{ when } T[i\dots] \neq P\]
This is essentially, ``given our choice of $q$, what is the probability we get a false positive?''  $t$ and $p$ are given by the input, and have nothing to do with the probability.\\
\textbf{Proof:} 
\begin{enumerate}
    \item Think of $P$ and $T$ as binary strings.  Define $S = T[i\dots]$, a substring of length $P$.
    \item $p \leq 2^N$ and $t_i \leq 2^N$, where $N = |P|$ (We have two symbols for the length of $P$ now.)  Basically the numbers represented in less than $N$ bits.  These are our un-\textit{mod}'ed numbers.
    \item Define a false positive (FP) when $t_i' = p'$ if and only if $(t_i - p) \equiv 0 \mod q$.  In other words, $q$ divides $|t_i - p|$.  Let $|t_i - p| = D$.  Now, we rephrase the question: given $D$, how many prime $q$'s can divide it evenly?
    \item $D \leq 2^N$.  This must be true because $D = |t_i - p|$, which both numbers must be smaller than $2^N$.  Now, we write $D = q_1 q_2 q_3 \dots q_k$, where $q_i \forall i < k$ are primes.  This is a prime factorization of $D$.  $k$ is the number of prime factors.  We know that $k \leq N$ since $D \leq 2^N$, and the prime factorization of $D$ is maximal when every factor is a $2$.  This places an upper bound on the number of numbers that divide $D$.
    \item The probability we picked a ``bad'' $q$ is the number of (bad choices / possible choices).  This fraction must be less than $\frac{N}{\text{number of primes less than }M} \leq M$.
    \item Pick $M$ so there are $5N$ primes less than or equal to $M$.  
    \[\Rightarrow \frac{\text{\# bad}}{\text{\# choices}} \leq \frac{N}{5N} = 0.2\]
    \item QED.
\end{enumerate}
The theorem allows us to pick an $M$, so we do so.  We could pick the theorem to be any arbitrary number by changing our decision of $M$.  However, we want $M$ to be big for the theorem, but also we cannot make it too big.  We now need a practical way to pick $M$.\\\\
\textbf{Fact:} If we want $n \geq 4$ primes between $1$ and $M$, it suffices to choose $M \geq 2n \log_2 n$.  We will not be proving this.
\begin{itemize}
	\item In our application, we need $5N$ primes less than $M$.  Therefore, we need $M \geq 2 \cdot (5N) \log_2 (5N) = 10N \log_2 (5N)$.  This is an upper bound on the size of $M$.
	\item This implies that we need $O(\log M)$ bits to represent our $q$.  Since $q \leq M$, all our numbers are bounded by $M$ bits.
	\item Plugging in, we get $O(\log(10 N \log (5N))) = O((\log 10N) + (\log \log 5N)) = O(\log N)$.  $N$ is the length of our pattern!
	\item This is definitely fast enough to be feasible for our use case of trying to fit all the computation in a register, because just to index $T$ (e.g., write down $T[i]$), at least $\log N$ bits are required.
\end{itemize}
The last thing we want to address is the size of the constant, $0.2$.  It is a \textit{constant}, but it is also large enough that $20\%$ false positives is kind of bad.  There are two ways to decrease the constant: increase $M$, or repeat picking a prime $R$ times instead of just once.
\begin{itemize}
	\item Increasing $M$ increases the constant in front of the space complexity, which might not be possible if we are trying to fit registers.
	\item However, we can roll more times at a constant time cost instead.
\end{itemize}

\subsection*{The Shift-And Algorithm}
Another algorithm for exact matching.  For this algorithm, we create a matrix $M = \{0, 1\}^{m \times n}$, where $m$ is the length of $P$ and $n$ is the length of $T$.  
\begin{itemize}
	\item The rule for filling out the matrix is $M[i, j] = 1$ if and only if $P[i] == T[j]$ and $P[1, i - 1] == T[\text{ending at }j - 1]$.
	\item The answer is the bottom row of the matrix.
	\item This is reminiscent of dynamic programming; the box at $M[i, j]$ depends on $M[i - 1, j - 1]$ and whether the characters in the string match.
	\item Unfortunately, a naive method of filling in this matrix is slow.  ($O(|T|^2)$).  We need a faster way.
\end{itemize}

\subsubsection*{Filling in the Matrix Fast}
If instead of filling out the cells one at a time, we filled it using vectorized operations instead, then we can make this run in $O(|T|)$.
\begin{itemize}
	\item Notice that each column in the matrix looks at the same ending character in $T$.
	\item Since every cell is based on it's up-left diagonal, we can do a Shift-And.  Essentially, take the column to the left, shift the entire column down and right by one, and AND the entire binary vector.
	\item We can now write the fill-out operation as
	\[M[\cdot, j] = U(T[j]) \texttt{ \& } (M[i, j - 1] >> 1)\]
    \begin{itemize}
	    \item where $U(x)$ is a $|P|$-long bit vector, such that $U(x)[i] == 1$ if and only if $P[i] == x$. 
	    \item Example: if $P = AACCTGA$, then $U("A") = 1100001$.
    \end{itemize}
\end{itemize}

\subsubsection*{Time Complexity}
Even when we vectorize the algorithm, does bitwise-AND run in constant time?
\begin{itemize}
	\item The answer depends on how big $P$ is.  If the entire binary vector is less than length 32 (e.g., $P$ is 32 characters long), then the entire vector may fit into a single register.
	\item This allows the algorithm to be hardware-optimized, and run in $O(n)$ time if $P$ is short enough.
	\item In practice, this algorithm runs as efficiently as (or even better than) the others of the conditions are met.
\end{itemize}

\section*{Approximate Matching}
$\ell$-mismatch problem: Given $T$, $P$, and $L$, find all places in $T$ where $T$ == $P$, up to $L$ disagreements.
\begin{itemize}
	\item $L$ is the budget of how many mismatches we get.
	\item $\ell$ is how many mismatches we have.
\end{itemize}

\subsection*{Variation of Shift-And for Approximate Matching}
Similar to the normal Shift-And, we also create a matrix, but in 3-dimensions.  Every $\ell$ receives its own matrix, where
\begin{itemize}
	\item $M^\ell[i, j] := 1$ if $P[1 \dots i] = T[\text{ending at }j]$ using $\leq \ell$ mismatches.
	\item This time, our answer would be the bottom row of the $L$-th matrix.  $M^L[|P|, \cdot]$
\end{itemize}
We have a similar recursion algorithm for dynamic programming:
\[M^\ell[j] = M^{\ell - 1}[j] \texttt{ OR } \left(M^\ell[j - 1] >> 1 \texttt{ AND } U(T[j])\right) \texttt{ OR } \left(M^{\ell - 1}[j - 1] >> 1\right)\]
\begin{itemize}
	\item There are three cases:
	\item The middle one is what we did before, just on one matrix of the $\ell$ matrices.
	\item The left one is the case where we use up a mismatch.
	\item The right one is the case where we add to our mismatch ``budget'' and use it on the next character.
\end{itemize}

\section*{Searching for Multiple Patterns}
The problem statement: find every occurance in $T$: $P_1, P_2, \dots, P_k$.
\begin{itemize}
	\item If we ran KMP invididually, (each string takes $O(|T| + |P|)$), then the total would be $O(|T| + |P_1| + |T| + |P_2| + \dots + |T| + |P_k|) = O(K|T| + \sum_i |P_i|)$
	\item The $\sum_i |P_i|$ part isn't a problem, but the $K|T|$ is.  We want to reduce the section of scanning over the input multiple times.
	\item However, there is a runtime caveat though, which is with multiple strings, the number of occurrences we print out at the end scales quadratically with the length of $T$.  This is unavoidable since patterns can overlap.
\end{itemize}

\subsection*{Aho-Corasick Algorithm}
We will define a keyword tree.  Start with $K(P_1, P_2, \dots, P_k)$.
\begin{itemize}
	\item For an example, let $P = \{\text{``abandon'', ``abduct'', ``abacus''}\}$.
\end{itemize}
\begin{center}
\begin{tikzpicture}
    \node[circle, draw] at (0, 0) (r) {};
    \node[circle, draw] at (0, -1) (a1) {};
    \node[circle, draw] at (0, -2) (b) {};
    
    \node[circle, draw] at (1, -3) (d1) {};
    \node[circle, draw] at (2, -4) (u1) {};
    \node[circle, draw] at (3, -5) (c) {};
    \node[circle, draw, fill=gray!50] at (4, -6) (t) {};

    \node[circle, draw] at (-1, -3) (a2) {};
    
    \node[circle, draw] at (-2, -4) (c1) {};
    \node[circle, draw] at (-3, -5) (u2) {};
    \node[circle, draw, fill=gray!50] at (-4, -6) (s) {};
    
    \node[circle, draw] at (0, -4) (n1) {};
    \node[circle, draw] at (1, -5) (d2) {};
    \node[circle, draw] at (2, -6) (o) {};
    \node[circle, draw, fill=gray!50] at (3, -7) (n2) {};

    \draw[->] (r) -- node[left] {a} (a1);
    \draw[->] (a1) -- node[left] {b} (b);
    \draw[->] (b) -- node[left] {a} (a2);
    \draw[->] (b) -- node[right] {d} (d1);
    \draw[->] (d1) -- node[right] {u} (u1);
    \draw[->] (u1) -- node[right] {c} (c);
    \draw[->] (c) -- node[right] {t} (t);
    
    \draw[->] (a2) -- node[left] {c} (c1);
    \draw[->] (c1) -- node[left] {u} (u2);
    \draw[->] (u2) -- node[left] {s} (s);

    \draw[->] (a2) -- node[right] {n} (n1);
    \draw[->] (n1) -- node[right] {d} (d2);
    \draw[->] (d2) -- node[right] {o} (o);
    \draw[->] (o) -- node[right] {n} (n2);
\end{tikzpicture}
\end{center}

\begin{itemize}
	\item Let $L(v)$ be the string spelled out from the root to $v$.
	\item Let $\ell_p(v)$ be the longest proper suffix of $L(v)$ that is a prefix of \textit{some} pattern
	\item Let $f(v)$ be the node representing $\ell_{p(v)}$
	\begin{itemize}
	    \item $f(v)$ is essentially a ``failure node'' or ``failure path'' that points backwards up the keyword tree.  Usage of this will be seen later.
    \end{itemize}
\end{itemize}

\subsubsection*{Using the Keyword Tree}
Starting at the root node, use $T$ as instructions as which branch of the tree to move down, similar to KMP.
\begin{itemize}
	\item If we reach the bottom node, then it is an accept state for one of the strings in $P$.
	\item If we ever reach a point where we cannot take a path, then we take the ``failure path'' to jump backwards up the tree to an earlier node.
	\begin{itemize}
	    \item Note that this is very similar to the \texttt{memo} array from KMP. In fact, if our keyword tree only had one pattern to match, it would be exactly the same.
    \end{itemize}
	\item Now, we just have to compute $f$.
\end{itemize}

\subsubsection*{But First, a Bug}
What happens if $P_i \subseteq P_j$?  That is, what if one pattern $P_i$ is a substring of another pattern $P_j$.  In this case, we could follow the path for $P_j$ and completely miss $P_i$.
\begin{itemize}
	\item To solve this problem, all we need to do is follow the failure function.  Or rather check it.  At any given point, if a pattern ends at the current character in $T$ we are traversing, the pattern would be a suffix.  The failure links have the purpose of matching suffixes to prefixes.
	\item However, there is a slight problem: The failure-links point to the \textit{longest suffix} that matches the prefix.  What if our pattern is not the longest?
	\begin{itemize}
	    \item In this case, we can follow the failure link again, and continue doing this until we reach the root.
    \end{itemize}
\end{itemize}
Now this introduces a new issue: if we have to follow all the failure links backwards to the root on each iteration, traversing them would take linear time\dots, for a total of quadratic time when iterating over the entire string.  
\begin{itemize}
	\item This is slow!
\end{itemize}
We can solve this problem by pre-processing the tree.  In addition to the failure links (F-links), we introduce output links (O-links), which act like F-links but skip all nodes that are not leaves.  In other words, they point from every node to all of the pattern matches.  Note that we can derive the O-links from the F-links list.
\begin{itemize}
	\item Now, we don't have to traverse the F-links on every iteration, we traverse the O-links instead.
	\item This might still take quadratic time in the worst case, but it is okay, because every node visited in the O-link represents an index that needs to be output.
	\item Our output, as established at the start, is un-improvably quadratic.  Therefore, traversing O-links do not take extra time.
\end{itemize}

\subsubsection*{Runtime}
This algorithm runs in $O(|T| + \sum_i |P_i| + z)$, where $|T|$ is the length of the string, $\sum_i |P_i|$ is the sum of the lengths of all the pattern strings, and $|z|$ is the size of our output. 
\begin{itemize}
	\item The $z$-term is quadratic by nature.
	\item This algorithm reads every character in our string $T$ once.  Building the keyword tree is in the $\sum_i |P_i|$ time, and traversing the tree is included in the $|z|$ term.
\end{itemize}

\subsubsection*{Computing the Failure Links, $f$}
Suppose we can compute the failure links in levels $1, \dots, f(v)$, where $v$ is a node somewhere in the tree, $u$ is the child of $v$, and $f(x)$ refers to the level of $x$.  Essentially, if we had F-links of all the nodes up to level $f(v)$, can we derive the F-links for node $u$ in constant time?
\begin{itemize}
	\item It turns out we can, and quite easily.  Consider the string segments $u$ and $v$ in $T$, where $u$ is one character longer than $v$.  By definition, the last $f(v)$ characters of $v$ match the first $f(v)$ characters of $v$.
	\item When we add a letter to the end of $v$ to form $u$, $x$ (the last character of $u$), we want to know the longest pattern in $v$ that is extendable by $x$.
	\item Notice that this is the purpose of the failure link of $v$!  By definition, the failure link of $v$ would give a prefix of $v$, so we can follow the failure links backwards until finding a pattern extendable by $x$.
\end{itemize}
\begin{center}
\textbf{Algorithm:} To compute $f(u)$, follow $f^*(\text{parent}(u))$ until you find a node extendable by label of edge $(v, u)$.
\end{center}
Now, the inductive hypothesis is complete; we only need the base case.  The base case happens to be the first character off of the root, in which case, the F-link of that node would point to the root.\\\\
What about the runtime?
\begin{itemize}
	\item We might have to follow multiple F-links from $v$ in order to find a node extendable by $x$.
\end{itemize}
Consider path $Q$ spelling pattern $P_i$.  We want to prove the total time to compute $f(u)$ for all $u$ on $Q$ is $O(|P_i|)$.
\begin{itemize}
	\item First, note that $|Q| = |P_i|$.
	\item Define $\ell_p(u) = L(f(u))$.  $\ell_p(u)$ is the string represented by $f(u)$.  $L(f(u))$ is the label of $f(u)$.  What we care about now is the length of the string, $|\ell_p(u)|$.  Since $u$ is parent of $v$, the length to $v$ is at most 1 larger.  That is, $\Delta |\ell_p| \leq 1$.  For simplicity, we will measure length in $\ell_p$ (and $\Delta |\ell_p|$) in dollars.
	\begin{itemize}
	    \item Going down the path gains \$1.  Following a F-link loses \$1.  Suppose we traverse the entire path normally.  If there are $q$ nodes in path $Q$, we gain at most $q$ dollars.  However, when we follow F-links backwards, we can never go past the root of the tree.  This implies the total amount of dollars at the end of a given path is \$0.
	    \item This in turn implies that the amortized amount spent on following all F-links equals the amount gained following the path $Q$.  Therefore, the total time to compute all F-links must be within $O(\sum_i |P_i|)$.
	    \item In implementation, we don't compute the patterns one at a time.  Rather, we compute all the patterns at the same time, following BFS order traversing the tree.  However, since all paths are amortized to their length, the entire tree would run in linear time.
    \end{itemize}
\end{itemize}



\end{document}